\let\negmedspace\undefined
\let\negthickspace\undefined
\documentclass[journal]{IEEEtran}
\usepackage[a5paper, margin=10mm, onecolumn]{geometry}
%\usepackage{lmodern} % Ensure lmodern is loaded for pdflatex
\usepackage{tfrupee} % Include tfrupee package

\setlength{\headheight}{1cm} % Set the height of the header box
\setlength{\headsep}{0mm}     % Set the distance between the header box and the top of the text

\usepackage{gvv-book}
\usepackage{gvv}
\usepackage{cite}
\usepackage{amsmath,amssymb,amsfonts,amsthm}
\usepackage{algorithmic}
\usepackage{graphicx}
\usepackage{textcomp}
\usepackage{xcolor}
\usepackage{txfonts}
\usepackage{listings}
\usepackage{enumitem}
\usepackage{mathtools}
\usepackage{gensymb}
\usepackage{comment}
\usepackage[breaklinks=true]{hyperref}
\usepackage{tkz-euclide} 
\usepackage{listings}
% \usepackage{gvv}                                        
\def\inputGnumericTable{}                                 
\usepackage[latin1]{inputenc}                                
\usepackage{color}                                            
\usepackage{array}                                            
\usepackage{longtable}                                       
\usepackage{calc}                                             
\usepackage{multirow}                                         
\usepackage{hhline}                                           
\usepackage{ifthen}
\usepackage{lscape}
\begin{document}

\bibliographystyle{IEEEtran}



\title{12.70}
\author{EE25BTECH11057 - Rushil Shanmukha Srinivas
}
% \maketitle
% \newpage
% \bigskip
{\let\newpage\relax\maketitle}

\renewcommand{\thefigure}{\theenumi}
\renewcommand{\thetable}{\theenumi}
\setlength{\intextsep}{10pt} % Space between text and floats

\numberwithin{equation}{enumi}
\numberwithin{figure}{enumi}
\renewcommand{\thetable}{\theenumi}
\textbf{Question: }$X$ bullocks and $Y$ tractors take 8 days to plough a field. If we halve the number of bullocks and double the number of tractors, it takes 5 days to plough the same field.
How many days will it take $X$ bullocks alone to plough the field?

\textbf{Solution: }

The total work done = W

Work done by 1 bullock in one day = x

Work done by 1 tractor in one day = y

Given,
\begin{align}
\myvec{8X&8Y}\myvec{x\\y} = W
\end{align}
\begin{align}
\myvec{\frac{5X}{2}&10Y}\myvec{x\\y} = W
\end{align}
The combined matrix is
\begin{align}
\myvec{8X&8Y\\
        \frac{5X}{2}&10Y}\myvec{x\\y} = \myvec{W\\W}
\end{align}
Forming the augmented matrix ,

\begin{align}
  \myaugvec{2}{8X & 8Y & W\\ \frac{5X}{2} & 10Y & W}
\end{align}
Using Gaussian Eliminations
\begin{align}
\myaugvec{2}{8X & 8Y & W\\ \frac{5X}{2} & 10Y & W} 
\xleftrightarrow{\;R_2 \to 4R_2 -5R_1}
 \myaugvec{2}{8X & 8Y & W\\-30X & 0 & -W} 
\end{align}
\begin{align}
30Xx = W
\end{align}
Number of days taken = 30
\end{document}